%*******************************************************
% Summary
%*******************************************************
\pdfbookmark[1]{Zusammenfassung}{summary}
\chapter*{Zusammenfassung}
\addcontentsline{toc}{chapter}{Zusammenfassung}
Diese Arbeit untersucht, inwieweit die Nachbearbeitung framebasierter Segmentierungen von Endoskopievideos durch probabilistische Sequenzmodelle verbessert werden kann. Ziel ist eine robuste Vorselektion nicht relevanter Bildbereiche, um rechenintensive Deep-Learning-Modelle zu entlasten und die Analysepipeline zu stabilisieren. Im Fokus stehen die Klassen \enquote{outside} und \enquote{low\_quality}, die für Anonymisierung und Qualitätskontrolle entscheidend sind.

Methodisch werden ein HMM und ein Conditional Random Field auf denselben regelbasierten Feature-Satz angewandt. Die Merkmale umfassen Farbstatistiken, Bewegungsinformationen und Texturgrößen, die vorab statistisch kalibriert und zu deterministischen Emissionsregeln verdichtet werden. Zur Erhöhung der Robustheit wird der ursprüngliche Labelraum auf eine reduzierte Menge semantischer Klassen abgebildet.

Die Ergebnisse zeigen, dass das Conditional Random Field insgesamt konsistenter als das HMM arbeitet und in der Gesamtleistung höhere Werte erreicht. Beide Modelle sind besonders stark in der Erkennung von \enquote{outside} und \enquote{low\_quality}, während die Erkennung pathologischer Befunde wie Polypen mit keiner der beiden Methodiken möglich ist. Die Laufzeitanalyse bestätigt zudem eine Echtzeitfähigkeit im Millisekundenbereich pro Frame.

Damit liefert die Arbeit einen praktikablen Filter für die Endoskopiesegmentierung, der Datenströme zuverlässig bereinigen kann und die Grundlage für eine nachgelagerte, präzisere Befunddetektion schafft.

\vfill
